\section{第一原则：阻尼本体论}

漂移线给出了一组位置，却没有说明这些位置凭什么成立。要让验证、准入、承责和定义真的接住价值，接受方必须能够确认一件事：给出承诺的那一方，是否有能力被后果改变。这个确认的前提比它看起来更基础：一个东西必须先有边界，才谈得上被改变。所以第一原则处理的是边界：边界如何被定义，代价如何证明边界的存在，以及意义为什么是一种需要持续投入才能维持的结构。

\subsection{边界的物理定义：交互即定义}

边界的形成方式，取决于它承担什么作用。容器的分隔由材料结构维持，细胞中的浓度差还要依靠代谢与转运；组织的权限边界则由身份、记录和程序确定。一个系统是否存在、边界能否被辨认、责任能否回到这个边界，分别需要相应证据。关于主体的自我报告可以被生成，因此判断还要考察身份是否连续、行为如何发生、后果能否追溯。

\subhead{边界是持续做功才能维持的差分}

细胞膜的物质分隔与膜两侧的电化学梯度，需要分别考察。离子泵持续消耗能量，补偿泄漏并维持相应梯度；泵送停止后，梯度按通道与浓度条件衰减，膜的材料结构却不因此立即消失。组织的准入、记录和追责也需要持续工作来维持功能。两者共同提示的是功能状态的维持问题，具体的材料机制与制度机制仍须分别说明。

把这段关系写成式子，就是差分在泵送停止之后的弛豫过程，

\begin{equation}
\frac{dB}{dt}=-\frac{B}{\tau}+P
\end{equation}

其中 $B$ 为边界的维持量，在离子泵这一情形里取膜两侧的离子浓度差，单位为摩尔每升；$P$ 为离子泵在单位时间内泵入的浓度差增量，单位为摩尔每升每秒；$\tau$ 为这个差分的弛豫时间，单位为秒，其大小决定了泵停止之后差分耗散到零所需的量级。泵一旦停止（$P=0$），方程的解为 $B(t)=B_0e^{-t/\tau}$，差分按指数衰减归零，正文所说的“几十分钟内耗散”对应的正是这个时间常数的量级；$P$ 恒定且为正时，系统停在 $B^{*}=\tau P$ 上，这个平衡值由泵送速率与弛豫时间共同决定。这个式子的含义是：边界不是一次形成的对象，而是一个必须由持续泵送对抗耗散才能维持的状态，输入撤掉之后它不是变弱，而是归零。适用条件需要一并写明：上式只刻画维持的结构，$B$ 取什么物理量、$\tau$ 与 $P$ 各为多少都取决于具体系统，正文只在离子泵这一情形下给出了时间尺度的量级。

\subhead{参照的必要性}

第二个支持来自参照的必要性。一个封闭系统内部的观测无法确定自身的绝对状态，这是相对性原理的直接后果：匀速直线运动中的观察者无法通过自身实验判断自己在运动，也无法判断自己的时钟是否在漂移。要确定状态，必须有外部参照。全球定位系统之所以需要在轨卫星与地面监测站持续比对，是因为星载原子钟的漂移率必须在与外部时间基准的比较中才能被标定；一台孤立的钟再精确，也无法自己发现自己慢了。这一点在工程上表现为反馈的不可替代：力控机械臂在接触之前无法判断抓取是否成功，因为它所缺的正是那个由对象给出的反作用力。系统对自己的认识，每一步都要借外界才能完成。

第三个支持来自生物学对环境边界的定义方式。Jakob von Uexk{\"u}ll 在二十世纪三十年代提出，动物并不生活在物理空间中，而生活在由它的感受与行动构成的周围世界里；蜱虫的周围世界几乎只包含丁酸的气味、体温和毛发触感三样信号，它对除此之外的一切不反应。\cite{uexkull1934streifzuge} 边界在这里不是几何尺寸，而是可交互项的集合。James J. Gibson 后来把这一点推进为可供性：环境被知觉到的不是中性的几何属性，而是一组可以直接支持行动的项，门把手被知觉为可拉的，台阶被知觉为可踩的。\cite{gibson1979} 控制论给出的版本更严格：调节器能够压制的扰动种类，取决于它自身的多样性能否覆盖这些扰动，缺少相应内部结构的调节器对这类扰动没有反应，也不会察觉自己被它改变。\cite{ashby1956cybernetics}

\subhead{边界用来选择通过什么}

把这三条放在一起，可以得到一个与直觉不同的结论：边界不是用来隔绝的，而是用来选择通过什么的。最硬的边界是晶体，它没有交换，也没有内部状态的可变性，因此它不会因为环境而改变，也不承担任何后果。活的边界相反，它必须让一部分东西通过，才能维持自己的差分。选择性通过的能力越强，边界对环境的依赖越深，它被改变的可能性也越大。完全封闭的边界意味着不能交换，不能交换意味着不能维持，所以封闭性不是边界的完成形态，而是它的死亡形态。

交互定义边界还有一层容易被忽略的性质：它要求可重复。一次碰撞只能证明那里有东西，不能证明那里有一个持续的主体。反复出现的反作用力才能画出轮廓，因为只有重复才能把随机扰动与稳定结构区分开。这也是为什么关系的边界、职业的边界和信誉的边界都需要时间来形成：它们由一系列可被检验的应答累积而成，而单次应答不构成轮廓。由此可以理解虚拟环境里的一个具体后果：当生成、推荐、安慰和执行都由系统代劳，人遇到的阻力被持续削减，可重复的反作用力被替换为可以随时修改的回应，边界的形成过程随之失去了材料。不是人的能力退化，是划定边界所需要的交互被取消了。

\subhead{边界是一个发育过程}

边界不仅在维持上是耗能的，在形成上还是一个发育过程。心理学里有一个可以直接观察的指标：把婴幼儿放在镜子前，先在他们鼻子上悄悄涂一个红点，如果他们伸手去摸自己的鼻子而不是镜中的影像，就说明他们已把这个影像识别为自己。这个反应通常在出生后十八到二十四个月之间出现，出现时间在个体之间差别很大，在此之前，婴儿对镜中影像作出的是对待另一个孩子的反应。\cite{amsterdam1972mirror} 这个转折点说明自我边界不是出生时就带有的设定，而是在与外部世界的反复交互中逐步建立起来的；它与前面三条支持同向：可以交互的项越多、越稳定，边界就越清楚，交互被替换为单向的回应时，边界就只能保持模糊。

免疫系统提供了另一个尺度上的例子。它区分自我与非我，靠的是发育过程中对自身抗原的耐受训练，以及此后持续的识别与清除；这套区分一旦失效，错误要么指向内部，要么放过外部入侵者，两个方向的后果都直接致命。\cite{burnet1959clonal} 它说明边界的准确性不是免费的：识别得越细，需要维持的机制越多，出错时的后果也越重。这也是边界问题在工程上总是表现为一组取舍的原因，把边界画宽，代价是放过；画窄，代价是误伤；而这两种代价由不同的人承担。

\subhead{这个原则的限度}

需要在这里留下限定，否则这个原则会被读成“阻力越大，边界越清楚”。一堵墙式的边界确实清楚，但它不交换，因此它不让任何东西发生。反过来，并非所有反作用力都构成边界：如果反作用力只落在同一侧（例如惩罚总是落在承受者身上，而施加者从不受影响），它划出的不是两侧的边界，而是一侧的单向通道。边界要求反作用力双向成立，也就是要求双方都能被对方的行动改变。这条限定解释了为什么代价必须落在能感受到它的位置上，也是下一节要处理的问题：一个东西凭什么证明自己确实付出了代价，以至于它的信号可以被依赖。

\subsection{昂贵信号：为什么必须付出代价才能证明}

\subhead{信号为什么必须昂贵}

孔雀尾羽等性状让信号成本成为演化研究的重要问题。Zahavi提出，高代价可能使信号与发送者状态联系起来；后续模型考察了成本、收益与类型如何共同维持诚实。\cite{zahavi1975mate,grafen1990biological} 性选择、直接指标、共同利益和约束也能形成可靠信号，不能把所有诚实都归为实际高消耗。

\begin{equation}
c(q_{\text{低}}) > b > c(q_{\text{高}})
\end{equation}

式中的 $c(q)$ 是质量为 $q$ 的个体发送这个信号的代价，$b$ 是信号被接受后带来的收益。不等式右半说明高质量个体发得起这个信号，左半说明低质量个体发不起，两个条件同时成立，信号才携带信息。这一结构并不限于生物。Michael Spence 在 1973 年用它分析劳动力市场：教育文凭之所以能向雇主传递能力信息，不是因为它教了什么，而是因为它对低能力者的成本更高。\cite{spence1973job} 累赘原理（handicap principle）本身并非没有争议，直到 2020 年仍有学者系统质疑它在经验证据上的牢固程度，认为它被接受得比证据所支持的更快。\cite{penn2020handicap} 本书沿用它的不是因为它在生物学上已经定论，而是因为它的形式结构足够清楚，在讨论信号成本时仍是一个有解释力的模型。把这三件事放在一起，可以得到一个与直觉相反的判断：\key{一个信号的价值不来自它表达了什么，而来自它不能被谁发出。} 如果所有人都能发出同一个信号，这个信号就不再区分任何人。

\subhead{三种不可伪造}

生成式 AI 改变的不是信号的数量，而是信号的伪造成本。过去要写出一篇结构完整的文章、画出一幅构图合理的图、写出一段能跑的代码，都需要多年训练；现在这些产物的伪造成本接近于零。伪造成本一旦下降，信号的信息量随之流失，这一点在文凭、作品集、简历这类传统信号载体上已经可以直接观察到。那么在这个背景下，什么仍然构成昂贵信号。前面那个不等式给出了判据的一半：代价与质量负相关；剩下的问题是，什么样的代价无法被低成本地复制。可以归结为三个条件。

第一是撤回条件。可以恢复的文件状态，仍可能消耗不可追回的电力和时间，也可能已影响他人。要检查的是所声称的投入与后果是否可追溯、恢复能达到什么范围。可回滚是一种工程优势，不会自动抹去全部投入。

第二是归属。工具可以代做一部分过程，人仍可检查、修改并承担最终判断。健身计划的文字可以外包，身体适应需要本人完成；道歉信可以辅助起草，关系修复还要靠双方行动。需要分别检查产出、学习、授权和承诺属于谁。

第三是形成与观察时间。有些能力需要重复、恢复和真实结果，信任需要反复兑现，某些安全性需要足够随访。更好的训练、设计与工具也能缩短准备。若缩短时间会改变产物，应指出对应机制与条件；单凭“三年压成三月”不能证明同一结果必然不可能。
\begin{equation}
t<\tau_{\mathrm{int}}\ \Longrightarrow\ \text{产物类型改变},
\qquad
\left.\frac{t}{\tau_{\mathrm{int}}}\right|_{\text{三年压到三个月}}=\frac{3}{36}=\frac{1}{12}
\label{eq:intrinsic-time}
\end{equation}
式中 $t$ 是实际占用的时长，$\tau_{\mathrm{int}}$ 是过程的内禀时间尺度，两者同为时间量纲，因此比值无量纲；$\tau_{\mathrm{int}}$ 按过程而定，技能养成取年计，信任建立取若干次兑现承诺的间隔，临床经验取积累到足够病例数所需的时间。判据说的是：只要 $t$ 落在 $\tau_{\mathrm{int}}$ 以下，输出的就不是同一种东西的劣质版本，而是另一种东西：三年训练压到三个月，得到的不是更快练成的技能，而是没有形成稳定结构的东西。辅助手段改变的是 $\tau_{\mathrm{int}}$ 前面的系数，也就是把同样的训练压缩到什么程度仍然落在门槛之上；跳过则是直接跨过门槛。时间刚性在这里表现为一个硬约束，而它与前两条的区别在于：外包改变的是能力的归属，加速改变的是产物还是不是同一种东西。

\subhead{高黏滞位置的几种投入}

把这三个条件应用到具体领域，可以看清价值正在向哪些位置移动，以及这些位置为什么会在加速环境中逆势变重。照护是其中最清楚的一例。一个人照顾失能的亲属，日复一日地处理起居、情绪和突发状况，这件事无法委托给一个更高效的执行者，因为被照顾的人需要的是这个特定的人在场。照护的时间不能被压缩，照护的失败不能被回滚，照护者的消耗不能被转移。正因为如此，照护一直难以被市场化定价，它的价值量级和它的报酬水平长期严重不匹配。这个错配本身就是信号：市场能定价的是可替换的服务，不能定价的是不可替换的在场。

科研在核心环节上有类似结构，但它的信号品种不同于照护。文献综述、数据整理、代码编写、初稿撰写，这些环节正在被大幅加速，真正留在原地的不是工作量，而是责任的落点。问题的选择、异常数据的判断、结论是否敢于承担的决断，仍然需要一个人押上自己的学术声誉。声誉之所以是昂贵信号，因为它不可回滚：一篇造假论文被揭露，损失无法通过撤回稿件来弥补。学术共同体用来认证的这种机制并不高效，它慢、保守、有时误伤，但它之所以还在运转，是因为没有更便宜的替代方案能产生同等可信度。

承诺和长期关系构成了第三个位置，它把时间刚性用到了极致。一段持续二十年的关系无法用二十年的对话记录来替代，因为关系的内容不是信息交换，而是双方在漫长时间里对彼此做出的、有代价的、被检验过的让步。这些让步在每一个时间点上都显得不划算，累积起来却形成了任何即时满足系统都无法复制的东西。不是它的内容更丰富，而是它的历史不能被重新生成，因为生成需要重来一遍那二十年。

身体劳动和公共信任代表了两类更基础的投入。身体劳动把人的存在直接接在物质世界上：耕作、建造、维修、手术、体育，这些活动的结果由物理定律给出，不由生成模型给出，一次用力不当造成的损伤不会因为事后修改了记录而消失。公共信任则是一种集体资产，它由无数次守约积累而成，由少数几次背约迅速摧毁。它的积累速率受制于人类兑现承诺的速度，不受制于信息传播的速度，这使它在加速环境中成为一个明显的慢变量。

\subhead{价值不是阻力本身}

到这里需要处理一个误解，这个误解在讨论这类问题时几乎必然出现：如果昂贵信号有价值，是不是阻力越大越有价值，痛苦越多越有价值。这个推论是错的，而且在历史上被反复用来为不必要的苦难辩护，所以必须把它和前面的论证切开。

信号理论和痛苦崇拜的区别在于，信号的价值来自它筛掉了什么，不来自它让谁痛苦。一个代价如果所有人都必须付出，它就不筛选任何人，因此不传递任何信息。历史上大量被强加的困难正是这种性质：它们消耗了人的一生，没有区分出任何东西，只证明了权力可以施加痛苦。这类代价不是昂贵信号，是纯粹的损耗。

由此可以给出一个比“阻力越大越好”精确得多的判据：不是代价的大小，而是代价与质量之间是否存在真实的负相关。一个有天赋的数学家在同一个训练上比一个没有天赋的人更轻松，这个训练就是信号；一个训练对所有人的难度完全一样，它就不是信号，只是门槛。这个区分让“哪些阻力应当保留”成为一个可以被分析的对象，而不是一个价值立场，因为“代价与质量是否负相关”是一个可以用证据来讨论的经验问题。

价值不是阻力本身，而是穿过必要阻力后形成的可靠结构。这句话可以反过来读：一个结构之所以可靠，是因为它经受过某种它无法回避的检验。检验不必是痛苦，但必须是真的：它必须能够拒绝，能够给出否定的结果，并且这个否定结果不能被撤销或者解释掉。

\subhead{伪造成本上升后的世界}

沿着这条线推下去，可以预期一个具体的变化方向：当符号生成的伪造成本接近于零，昂贵信号的供给会集中到那些伪造成本无法降低的位置上。这个转移不会自动发生，因为社会对信号的需求不会因为伪造变容易就消失，它只会先转移到更容易被伪造的替代品上。这与一个常见的观察吻合：在文凭贬值的同时，学历要求反而在上升，因为雇主仍然需要筛选，而手边可用的筛子只剩这一把。

真正的转向发生在需求端，而不是供给端。当人们逐渐发现大部分可生成的信号不再有信息量，他们会开始要求那些无法生成的证明：现场经验、可追溯的责任记录、可验证的失败史、可以联系的推荐人、被时间检验过的作品。这些要求的共同点是它们都不能被快速制造，因此在高效率环境中显得低效和麻烦，也因此在组织内部往往是最先被提议删掉的部分。但它们承担的功能，正是让判断有一个可以落脚的支点。

哪些阻力属于这一类，哪些只是应该被消除的浪费，是把上述判断转成行动时绕不开的问题。这个问题之所以紧迫，是因为在加速环境中，两类阻力在外观上难以区分：它们都表现为变慢、变麻烦、变得不划算，而它们的后果却方向相反。

代价与质量之间的这种关系还有另一面。付出代价不只是筛掉了谁，它还在系统内部留下了结构：训练在身体上留下能力，承诺在关系里留下记录，失败在判断里留下校准。这些结构不是一次性的产物，它们需要被持续维持。于是问题从“代价证明了什么”转向“代价留下了什么，以及留下来之后靠什么维持”。

\subsection{意义热力学：负熵构建}

\subhead{位移不等于做功}

力学里有一个区分，用来处理看起来动了、实际上没有能量转移的情形。当一个系统发生虚位移时，约束力所做的功为零；判断功的大小，看的不是位移本身，而是力与真实位移在同一方向上的乘积。这个区分可以搬来分析生成式系统的产出。模型可以在几秒之内生成一部史诗的框架、一幅构图完整的画、一份结构清晰的方案，屏幕上发生的位移非常可观，但在整个过程中被克服的物理阻力接近于零：没有痛感，没有风险，没有生命消耗，也没有任何一方被这个过程改变。位移发生了，做功没有发生。把这种情形称为虚功幻觉并不过分，因为它最容易被误读的部分恰恰是它的规模：产出越多，越像是有某种力量被付出了。

意义不服从位移的判据。一个人用十年写出一部书，和模型用十秒生成一部结构相似的书，两者在文本长度上可以接近，在可检索性上甚至后者更高，但它们在现实中的存在方式不同。前者的十年是真实消耗掉的时间，作者在这段时间里被写作改变，他的判断、关系和身体都留下了痕迹，这些痕迹构成书的分量；后者没有经历任何这样的过程，它是一次位移，不是一次做功。差别并不神秘，它可以被还原为一个物理问题：有没有一个系统因为这件事而不可逆地改变了状态，并且这个改变必须被维持。

\subhead{负熵结构的维持成本}

把意义的构建放到热力学里看，它对应的是负熵结构的形成与维持。第二章已经交代过薛定谔的负熵论证和普里高津的耗散结构\cite{schrodinger1944,nicolis1977,prigogine1977nobel}：这类有序结构依赖持续的通量和非线性相互作用，一旦输入输出停止，结构就瓦解、冷却或回到更接近平衡的状态。这里要用的是它的一个推论：结构的形成成本与维持成本是两笔账，而且后者通常更长。

这两笔账的分工可以写成熵平衡的形式，

\begin{equation}
\frac{dS_{\text{系}}}{dt}=\dot S_{\text{产}}+\dot S_{\text{流}},\qquad \dot S_{\text{产}}\ge 0
\end{equation}

$S_{\text{系}}$ 为被维持的结构自身的熵，单位为焦每开；$\dot S_{\text{产}}$ 为结构内部不可逆过程产生的熵，单位为焦每开每秒，按热力学第二定律它恒不为负；$\dot S_{\text{流}}$ 为结构与环境交换带来的净熵流，单位为焦每开每秒。要维持 $\frac{dS_{\text{系}}}{dt}\le 0$ 的低熵状态，必须有 $\dot S_{\text{流}}\le-\dot S_{\text{产}}<0$，也就是持续的负熵输入，而它在环境一侧对应着等量的熵增。这个式子的含义是：形成只改变初始条件，维持则要求这条熵流一直流下去，流速停止时结构不是静止，而是瓦解，把这一条写成式子，就是下面这句话的形式依据：生成降的是进入门槛，维持靠的是持续通量。适用的边界必须写清楚：这是第二定律对开放系统的常规表述，它只说明低熵状态的维持需要一个持续的通量，不说明这个状态是否值得维持，也不构成任何价值判断，把它读成价值证明正是本节点名的第三种误用。

这两笔账在意义结构上分得很开。一个想法被提出的成本可以很低，数学史上不少关键构造是在很短的时间里被写下来的；但一个想法要成为共同体的知识，需要被反复检验、引用、教学和修订，这个过程以几十年计。一部经典文本的存续同样如此：文本本身不会自己保存下去，它靠持续的抄写、注释、翻译、讲授和再版活着。中世纪欧洲的许多古典著作之所以能留下来，不是因为它们写得好，而是因为有一代又一代人在羊皮纸上重抄它们，抄写本身就是维持成本的支付。一旦这条维护链断掉，结构不会静止在那里等待被重新发现，它会退化为不可读的噪声。

人体提供了更直接的量级。人脑约占体重的 2\%，却消耗静息状态下全身氧耗的约 20\%，是身体里最昂贵的器官之一。\cite{raichle2002} 这一笔开销买到的不是即时的产出，而是维持一个能够持续重组自身连接的结构：记忆的巩固依赖反复的提取与睡眠期间的重放，技能的稳定依赖成千上万次重复，判断的校准依赖真实后果的多次反馈。这些过程的共同点是它们的速率由生理过程决定，不因为外部信息处理变快而变快。把三年的训练压缩到三个月，得到的不是更高效的训练，而是另一种东西。这一点在信号那一节已经出现过，这里显示的是它的热力学来源。

生成成本与维持成本需要分别估计。更好的设计和工具可以同时降低两者，结构的存续则仍取决于实际维护、使用和传承。计划若只按候选生成的速度安排，而没有配置相应的校验与维护能力，差额便可能以不可读、不可信或失修的形式出现。

\subhead{这个框架的限度}

持续投入说明一个结构如何被维持，价值判断还要考察它维持了什么。会议、流程和制度都可能成本很高，所获得的结果却不足以支持这些支出。需要保留的功能应当能够被指出、被检查，并与其他实现方式比较。成本高低本身不给出取舍。

第三种误用是把负熵当成价值证明。它的形式很固定：把“需要维持”读成“值得维持”，中间跳过一步，而跳过的那一步恰是关键。两样东西可以同样需要持续投入、同样要付出代价才不退回噪声，一个值得保留，一个应当尽早结束，投入的账上看不出区别；要区分它们，靠的是这套框架之外的判据，而给出判据的人要承担判错的后果。真正能从这套框架里拿走的是一条结构性的结论：任何被反复识别的意义都对应一条持续的投入链，而这条链上的位置不能由生成能力替代，只能由愿意承担维持成本的主体占据。

维持成本还有一个分布问题。谁支付维持成本、谁从结构的存在中获益，这两个问题在意义上和在物质上一样不能回避。一项被广泛使用的标准、一种被普遍接受的规范、一门被代代传授的手艺，它们的维护者往往是具体的人，而受益者却是整个网络。这种不对称会使维护投入长期低于社会最优水平，除非有制度性的安排来承担这笔开销。到这里，问题不再是意义需要多大的投入，而是投入的方向由谁决定：谁有权决定维持什么、放弃什么，以及凭什么。这正是下一节要处理的内容。
